% Drop-in replacement for \subsection{Binary system} of the SPICE paper (mnras_spice_draft_v2.tex,
% from "\subsection{Binary system}" to the end of fig:tz_fornacis_eclipse).
% 2026-11-02: a pure literature reproduction -- NOTHING is fitted. Gallenne et al. (2016) parameters and light-curve
% geometry, Andersen et al. (1991) photometric ephemeris, phase-3 emulator, exact per-element synthesis.
% Labels unchanged (sec:binary_system, tab:tz_fornacis_parameter_values, fig:tz_fornacis_lc, fig:tz_fornacis_spectra,
% fig:tz_fornacis_eclipse); new label tab:tz_fornacis_models.
% Figures: tz_fornacis/tzfor_literature_figures.ipynb on out/tzfor_pure_p3_pure.npz
%   paper_plots/tz_fornacis/tzfor_lit_lc.pdf -> figures/tz_for_lc.pdf; tzfor_lit_spectra.pdf -> figures/tz_for_spectra.pdf;
%   tzfor_lit_lc_eclipses.pdf -> figures/tz_for_lc_eclipses.pdf
% The parameter table's second column is a wrapping p{} column (\usepackage{array}).
% Numbers (Gadi /scratch/y89/mj8805/litrepro/): pure/pure_summary.txt (SPICE, PHOEBE), pure/port_summary.txt (PHOEBE model
% in SPICE), interp/out/lit9_summary.txt (sensitivity tests), interp/out/fe1_ew (equivalent widths).
% Scoring of every light curve (SPICE and PHOEBE alike): tz_fornacis/tzfor_phot_rescore.py --ephemeris 2445032.6091 75.66675
% --dphi 0 --nightly none. Bib: tzfor_paper_section.bib (Andersen1991, Mayor2003, Rozanski2022; Czekala2015 only for the
% optional paragraph at the end of this file).

\subsection{Binary system}
\label{sec:binary_system}

To test the binary-star modelling of \spice\ against real data we selected a wide, fully detached
eclipsing system whose components fall inside the domain of the emulator's training grid and for which
both photometric and spectroscopic time series are public: TZ~Fornacis, a K giant with an F-type
companion in a 75.7-day circular orbit. The data are the Str\"omgren $b$ and $y$ light curves of
\citet{Clausen1991} (1169 epochs per band, differential photometry against HR~970) and 21 HARPS
\citep{Mayor2003} spectra, the pipeline-reduced products of the ESO Science Archive Facility obtained
between 2009 and 2015, eleven of them under the Araucaria programme 190.D-0237 whose radial velocities set
the orbit of \citet{tz_fornacis_orbit}; we continuum-normalised them with \texttt{SuppNet} \citep{suppnet}.

We do not fit the system: we ask whether its published solution, put through \spice\ unchanged,
reproduces both data sets. Every input is taken from the literature
(Table~\ref{tab:tz_fornacis_parameter_values}). From \citet{tz_fornacis_orbit} come the masses, the orbit,
the temperatures they adopt (4930 and 6350\,K), their spectroscopic $[\mathrm{Fe/H}]$ ($+0.02$ and
$-0.05$), microturbulence (0.98 and 3.10\,km\,s$^{-1}$), projected rotation (6.1 and 45.7\,km\,s$^{-1}$)
and the separate systemic velocities they fit to the two stars (17.99 and 18.35\,km\,s$^{-1}$). The
geometry is that of their light-curve model, from which the adopted secondary temperature was derived:
radii of 8.34 and 3.97\,$R_\odot$ (the fractional radii of \citealt{Andersen1991} on their semi-major
axis), Roche surfaces with each star rotating at $v\sin i / \sin i$, so that the fast-rotating secondary
is flattened by about 1 per cent, and Lucy gravity darkening. Their projected rotation comes from
broadening functions measured against templates without macroturbulence and therefore already contains
it, so no macroturbulence is added. The photometry is phased with the ephemeris that
\citet{Andersen1991} derived from these same light curves. The surfaces are \phoebe\ meshes imported as
described in Section~\ref{sec:phoebe_integration}, the specific intensities come from the \tpayne\
emulator of Appendix~\ref{sec:appendix_transformer_emulators}, and every spectrum and light curve is the
exact per-element sum of Section~\ref{sec:integrating_flux_from_mesh}, with no pixels masked. Nothing is adjusted to the data except one zero point per photometric band, which the
differential photometry leaves undetermined; there is no phase offset, no per-night zero point and no
correction of the spectroscopic continuum. Of the 2338 photometric points, 7 lie more than 15\,mmag from
the model and are not counted.

The light curves are reproduced to the scatter of the photometry
(Figs.~\ref{fig:tz_fornacis_lc} and \ref{fig:tz_fornacis_eclipse}): $\chi^2/N = 0.96$, with residuals of
3.9 and 3.6\,mmag in $b$ and $y$, smaller inside the eclipses (3.5\,mmag in both bands) than outside them.
The deep eclipse is centred on the published ephemeris to 0.04\,h; the shallow one falls 0.2--0.4\,h late,
the displacement expected from the $e\cos\omega = 0.0002 \pm 0.0001$ that \citet{Andersen1991} measured in
these light curves before adopting a circular orbit, which we keep.
\phoebe\ evaluated at the same parameters and scored in the same way gives $\chi^2/N = 0.93$ with its
\texttt{ck2004} atmospheres and irradiation, 0.95 without irradiation, and 1.14 with its \texttt{phoenix}
atmospheres (Table~\ref{tab:tz_fornacis_models}). To separate the radiative transfer from the surface
model we also ported the \phoebe\ model itself into \spice, including its per-element temperatures, so
that only the specific intensities differ: this gives $\chi^2/N = 0.951$, against 0.949 for \phoebe\
without irradiation and 0.960 for \spice's own gravity-darkening law. The two codes therefore agree to the
precision of the photometry, and the remaining difference to the full \phoebe\ model, 0.02 in
$\chi^2/N$, is the light reflected between the stars (under 1\,mmag), which \spice\ does not compute;
\phoebe's irradiation raises the surface temperatures by less than 3\,K. With the interferometric radii of
\citet{tz_fornacis_orbit} (8.28 and 3.94\,$R_\odot$) instead of their light-curve radii (and the secondary
rotating at 38\,km\,s$^{-1}$, as in our earlier meshes), $\chi^2/N$ rises to 1.14, and with their Fe\,\textsc{i} equivalent-width temperature for the secondary, 6650\,K, it rises
by a further 0.8, with the eclipse depths visibly wrong: the light curves confirm the adopted 6350\,K, as
\citet{tz_fornacis_orbit} concluded from their $H$-band flux ratio.

The spectra are reproduced to $\chi^2/N = 17.5$ against the \texttt{SuppNet}-normalised data as they are,
and to 12.7 if a smooth continuum correction is allowed for each spectrum (cubic splines with 40\,\AA\
knots), so a quarter of the mismatch is in the normalisation of the observations.
Figure~\ref{fig:tz_fornacis_spectra} shows how the blend is built at three orbital phases. The two stars
separate cleanly because the giant's narrow lines and the secondary's rotationally broadened ones are
distinguishable at every phase, and their dilution follows from the light ratio rather than being a free
parameter. The line mismatch that remains has three identified sources. First, the spectral grid has its
own metallicity scale: analysed with the same pipeline, HARPS spectra of Gaia FGK benchmark giants and
the solar flux atlas return $[\mathrm{Fe/H}]$ 0.21--0.26\,dex below their reference values, and lowering
the literature $[\mathrm{Fe/H}]$ of both stars by 0.2\,dex reduces $\chi^2/N$ by 2.3 (with the continuum
correction). Second, the microturbulence scale differs from that of the ATLAS9/\texttt{MOOG} analysis of
\citet{tz_fornacis_orbit}. Repeating their Fe\,\textsc{i} equivalent-width analysis with our models
(about 200 isolated lines measured identically in the observed and the synthetic composite spectra)
recovers their metallicity at their microturbulence, a median of $+0.13$, but leaves a trend of the line
abundances with reduced equivalent width; the trend vanishes at $\xi_1 \approx 1.9$--2.0\,km\,s$^{-1}$,
where the lines give $[\mathrm{Fe/H}] \approx -0.32$ to $-0.37$. The two analyses reproduce the same
equivalent widths with different $(\xi, [\mathrm{Fe/H}])$ pairs, so the published pair is not the best one
for our grid. Third, the residuals are correlated along the spectrum over 33 pixels (27\,km\,s$^{-1}$, the
width of a line), so $\chi^2/N$ is not a goodness-of-fit statistic here and only its differences between
models are meaningful. Adding a macroturbulence of 5 and 6\,km\,s$^{-1}$ on top of the published
$v\sin i$ lowers it by 1.1, a hint that the broadening-function $v\sin i$ underestimates the total line
width for our profiles.

As a check of \spice's own occlusion, the light curves computed from two \spice\ icospheres on the
Keplerian orbit agree with those from the imported \phoebe\ meshes to 0.6\,mmag at the same parameters.
The zoomed-in view of the eclipses in Figure~\ref{fig:tz_fornacis_eclipse} also shows the expected
wavelength dependence: with limb darkening and the partial occultation of the primary modelled, the
eclipse profile is steeper in bluer filters. This dependence is more commonly emphasised for planetary
transits, and its implications for the use of \spice\ in exoplanet applications will be explored in
future work.

\begin{table}
    \centering
    \caption{The TZ Fornacis model: the published solution, used unchanged. Stellar and orbital parameters
    from \citet{tz_fornacis_orbit}; radii from their light-curve model; ephemeris of the photometry from
    \citet{Andersen1991}.}
    \begin{tabular}{l>{\raggedright\arraybackslash}p{0.55\columnwidth}}
        \hline
        \textbf{Parameter} & \textbf{Value} \\
        \hline
        Masses & 2.057, 1.958\,$M_\odot$ \\
        Radii & 8.34, 3.97\,$R_\odot$ \\
        $T_{\rm eff}$ & 4930, 6350\,K \\
        $\log g$ & 2.91, 3.52--3.54 (Roche surfaces) \\
        $[\mathrm{Fe}/\mathrm{H}]$ & $+0.02$, $-0.05$ \\
        Microturbulence $\xi$ & 0.98, 3.10\,km\,s$^{-1}$ \\
        $v \sin i$ & 6.1, 45.7\,km\,s$^{-1}$ \\
        Systemic velocity $\gamma$ & 17.99, 18.35\,km\,s$^{-1}$ \\
        Macroturbulence & none (contained in $v \sin i$) \\
        Gravity darkening & $\beta_{\rm bol} = 0.32$ (convective envelopes) \\
        Period & $75.66647$\,d \\
        Eccentricity & 0 \\
        Inclination & $85.68^\circ$ \\
        Argument of periastron & $269.93^\circ$ \\
        Longitude of the ascending node & $65.99^\circ$ \\
        Photometric ephemeris & Min\,I = HJD\,2445032.6091 + 75.66675\,$E$ \\
        \hline
    \end{tabular}
    \label{tab:tz_fornacis_parameter_values}
\end{table}

\begin{table}
    \centering
    \caption{TZ Fornacis at the published parameters, with nothing fitted: reduced $\chi^2$ of the Str\"omgren
    photometry (2331 points; one zero point per band) with the $b$ and $y$ residual rms, and of the HARPS
    spectra ($3.0\times10^{6}$ pixels; in parentheses with a smooth continuum correction per spectrum). The
    spectral $\chi^2$ is inflated by correlated residuals; only its differences are meaningful.}
    \begin{tabular}{lccc}
        \hline
        Model & $\chi^2_{\rm phot}/N$ & rms $b$, $y$ [mmag] & $\chi^2_{\rm spec}/N$ \\
        \hline
        \spice                                              & 0.960 & 3.9, 3.6 & 17.5 (12.7) \\
        \spice, \phoebe's surface temperatures              & 0.951 & 3.8, 3.6 & 17.5 (12.7) \\
        \phoebe, \texttt{ck2004}                            & 0.927 & 3.8, 3.5 & -- \\
        \phoebe, \texttt{ck2004}, no irradiation            & 0.949 & 3.9, 3.5 & -- \\
        \phoebe, \texttt{phoenix}                           & 1.135 & 4.2, 3.9 & -- \\
        \hline
    \end{tabular}
    \label{tab:tz_fornacis_models}
\end{table}

\begin{figure*}
    \centering
    \includegraphics[width=\textwidth]{figures/tz_for_lc.pdf}
    \caption{Str\"omgren $b$ and $y$ light curves of TZ Fornacis from \citet{Clausen1991} against \spice\ at
    the published parameters (Section~\ref{sec:binary_system}), with the residuals beneath. The model is the
    exact per-element synthesis through the imported \phoebe\ meshes, phased with the ephemeris of
    \citet{Andersen1991}; nothing is fitted except one zero point per band. Crosses are the 7 points more
    than 15\,mmag from the model.}
    \label{fig:tz_fornacis_lc}
\end{figure*}

\begin{figure*}
    \centering
    \includegraphics[width=\linewidth]{figures/tz_for_spectra.pdf}
    \caption{How the blended spectrum of TZ Fornacis is built in the literature model of
    Section~\ref{sec:binary_system}, in the Mg\,b region. Columns are three orbital phases along the
    radial-velocity cycle: maximum separation with the giant receding, near conjunction with the two line
    systems blended, and maximum separation with the giant approaching; the sketches show the
    configuration, with separations to scale and the discs enlarged. Rows are the giant's contribution,
    the secondary's contribution, and their sum against the HARPS spectrum with the residual. Nothing is
    fitted, including the continuum: the offset of a few per cent between the observed and the model
    continuum is that of the \texttt{SuppNet} normalisation.}
    \label{fig:tz_fornacis_spectra}
\end{figure*}

\begin{figure*}
    \centering
    \includegraphics[width=\textwidth]{figures/tz_for_lc_eclipses.pdf}
    \caption{The two eclipses of TZ Fornacis in Str\"omgren $b$ and $y$, the same model and data as
    Figure~\ref{fig:tz_fornacis_lc}. The deeper eclipse is the secondary passing behind the giant, the
    shallower its transit.}
    \label{fig:tz_fornacis_eclipse}
\end{figure*}

% ---------------------------------------------------------------------------
% OPTIONAL paragraph (fits; not part of the literature reproduction). If wanted, insert after the spectra paragraph:
%
% For comparison, fitting the temperatures, metallicities, macroturbulence, the giant's microturbulence and its C, N
% and O to both data sets (interferometric-radius geometry fixed) reaches $\chi^2_{\rm spec}/N = 8.8$ (with the
% continuum correction) at $T_{\rm eff} = 4869$ and 6350\,K, $[\mathrm{Fe/H}] = -0.27$ and $-0.30$, $\xi_1 =
% 1.5$\,km\,s$^{-1}$, $[\mathrm{C/Fe}] = -0.26$ and $[\mathrm{N/Fe}] = +0.35$, independent of the relative weight of the
% two data sets (to 10\,K for spectroscopic weights from 1 to $3\times10^{-3}$); a Gaussian-process model of the
% correlated residuals \citep[as in][]{Czekala2015} moves $T_{\rm eff,2}$ to 6429\,K. The fit recovers the adopted
% temperatures and lands on the metallicity and microturbulence scales identified above.
% (Gadi: interp/ai-fit-tz-vtw*, ai-fit-tz-gp; photometry of these fits used per-night zero points.)
% ---------------------------------------------------------------------------
% For \section*{Acknowledgements} (ESO's required form of words):
%
% Based on observations collected at the European Southern Observatory under ESO
% programmes 082.D-0499, 083.D-0549, 084.D-0591, 085.C-0614, 086.D-0078,
% 094.D-0056 and 190.D-0237, obtained from the ESO Science Archive Facility.
%
% For \section*{Data Availability}, add after the Clausen citation:
%
% The HARPS spectra of TZ Fornacis are public in the ESO Science Archive Facility
% under the programme identifiers listed in the Acknowledgements.
% ---------------------------------------------------------------------------
